
\begin{frame}[fragile]{SQL, agrégation}

Cette session présente les agrégats en SQL. Elle consiste à regrouper des lignes
et à appliquer à chaque groupe une fonction d'agrégation.

\vfill

Contenu:

\begin{redbullet}
  \item La clause \sqlkw{group by}
  \item Fonctions d'agrégation
  \item La clause \sqlkw{having}
\end{redbullet}

\vfill
\begin{blocgauche}
\alert{Ces diapositives correspondent au support en ligne disponible sur
le site \url{http://sql.bdpedia.fr/}}
\end{blocgauche}

\end{frame}

%%%%%%ù
\begin{frame}[fragile]{Principe général}

On définit des \alert{groupes} de lignes ayant en commun une ou plusieurs valeurs.

\vfill

On ramène chaque groupe à une seule valeur en appliquant une \alert{fonction d'agrégation}.

\vfill

Cas le plus simple: un seul groupe, obtenu par un bloc \texttt{select from where}.


\begin{minted}[baselinestretch=.9,
               fontsize=\small,
               framesep=2mm]{sql}
select count(*) as nbPersonnes, count(prénom) as nbPrénoms, 
         count(nom) as nbNoms
from Personne
\end{minted}

\vfill

\begin{tabular}{c|c|c}
\textbf{nbPersonnes} & \textbf{nbPrénoms} & \textbf{nbNoms}\\ \hline 
7 & 6 & 7\\ \hline 
\end{tabular}

\end{frame}

%%%%%%ù
\begin{frame}[fragile]{Le \texttt{group by}}

La clause \sqlkw{group by} \texttt{att1, ..., attn} \alert{partitionne} le résultat d'un bloc
\sqlkw{select from where} en fonction des \texttt{att1, ..., attn}

\vfill 

Chaque \alert{groupe} contient les lignes qui partagent les mêmes valeurs pour \texttt{att1, ..., attn}.

\vfill

\begin{minted}[baselinestretch=.9,
               fontsize=\small,
               framesep=2mm]{sql}
select  idAppart, sum(quotePart) as totalQP
from    Propriétaire
group by idAppart
\end{minted}

\vfill
\begin{blocgauche}
Procède en deux étapes: d'abord on groupe, puis on agrège.
\end{blocgauche}
\end{frame}


%%%%%%ù
\begin{frame}[fragile]{Décomposons\,: l'étape de regroupement}

On obtient une structure intermédiaire, avec autant de lignes que de valeurs
distinctes pour les attributs de regroupement (ici, \texttt{idAppart}).

\vfill

\begin{tabular}{c|c}
\textbf{idAppart} & Groupe (\textbf{idPersonne}, \textbf{idAppart}, \textbf{quotePart}) \\ \hline 
100 & \begin{tabular}{c} 
            [(1, 100, 33), (5, 100, 67)] \\ 
       \end{tabular}                
       \\ \hline 
101 &  \begin{tabular}{c} 
            [(1, 101, 100)] \\ 
       \end{tabular}                \\ \hline 
102 &  \begin{tabular}{c} 
            [(5, 102, 100)] \\ 
       \end{tabular}                \\ \hline 
103 &  \begin{tabular}{c} 
            [(2, 103, 100)] \\ 
       \end{tabular}                \\ \hline 
104 &  \begin{tabular}{c} 
            [(2, 104, 100)] \\ 
       \end{tabular}                \\ \hline 
201 &  \begin{tabular}{c} 
            [(5, 201, 100)] \\ 
       \end{tabular}                \\ \hline 
202 &  \begin{tabular}{c} 
            [(1, 202, 100)] \\ 
       \end{tabular}                \\ \hline 
\end{tabular}

\vfill
\begin{blocgauche}
\alert{Ce n'est pas une table en première forme normale.}
\end{blocgauche}
\end{frame}

%%%%%%ù
\begin{frame}[fragile]{L'étape d'agrégation}

La fonction d'agrégation ramène un groupe à une valeur
\vfill

\begin{tabular}{c|c}
\textbf{idAppart} & SUM(\textbf{quotePart}) \\ \hline 
100 & SUM (33, 67) = 100 \\ \hline 
101 &  SUM (100) = 100\\ \hline 
102 &  SUM (100) = 100\\ \hline  
103 &  SUM (100) = 100\\ \hline 
104 &  SUM (100) = 100\\ \hline  
201 &  SUM (100) = 100\\ \hline  
202 & SUM (100) = 100\\ \hline  
\end{tabular}

\vfill
\begin{blocgauche}
\alert{Cette fois c'est une table en première forme normale.}
\end{blocgauche}
\end{frame}


%%%%%%ù
\begin{frame}[fragile]{La clause \sqlkw{having}}

Exprime un critère de sélection sur le résultats de la fonction d'agrégation.
\vfill

\alert{Bien distinguer de la clause \sqlkw{where} qui s'applique aux nuplets}

\vfill

\begin{minted}[baselinestretch=.9,
               fontsize=\small,
               framesep=2mm]{sql}
     select  idAppart, count(*) as nbProprios
     from    Propriétaire
     group by idAppart
     having count(*) >= 2
\end{minted}

\vfill

\begin{tabular}{c|c}
\textbf{idAppart} & \textbf{nbProprios}\\ \hline 
100 & 2\\ \hline 
\end{tabular}

\end{frame}

%%%%%%%%%
\begin{frame}{À retenir}

Agrégats = extension de SQL.

\vfill

\begin{redbullet}
  \item  S'applique au \alert{résultat} d'une requête standard
  \item Partitionne en groupes de nuplets partageant les mêmes valeurs de regroupement
  \item Réduit chaque groupe à une valeur grâce à une fonction d'agrégation
  \item On peut filtrer les groupes obtenus avec \sqlkw{having}
\end{redbullet}

\vfill

\begin{blocgauche}
Attention aux valeurs  à \texttt{null} et aux doublons.
\end{blocgauche}
\end{frame}

